\begin{Answer}{6.6}
(a) Since we assumed that $n\geq 1$, $-3n$ is certainly negative.  In other words,
$-3n\leq 0$.  That's why in the first step we could say that $5n^2-3n+20\leq 5n^2+20$.
(b) We used the fact that $20\leq 20n^2$ whenever $n\geq 1$.
If either of these solutions is not clear to you,
you need to brush up on your algebra.
\end{Answer}
\begin{Answer}{6.7}
This is incorrect.
It is {\em not true} that $-12n\leq -12n^2$ when $n\geq 1$.
(If this isn't clear to you after thinking about it for a few minutes, you may need to
do some algebra review.)
In fact, that error led to the statement
$4 n^{2} - 12n + 10\leq 2n^2$ which cannot possibly be true as $n$ gets larger
since it would require that $2n^2-12n+10\leq 0$.
This is not true as $n$ gets larger.
In fact, when $n=10$, for instance, it is clearly not true.
But it {\em is true} that $-12n<0$ when $n\geq 0$, so instead of replacing it
with $-12n^2$, it should be replaced with $0$ as in previous examples.
\end{Answer}
\begin{Answer}{6.8}
(a) Sure.  Add the final step of $25n^2\leq 50n^2$ to the algebra in the proof.
In fact, any number above $25$ can easily be used.
Some values under 25 can also be used, but they would require a modification
of the algebra used in the proof.
The bottom line is that there is generally no `right' value to
use for $c$.  If you find a value that works, then it's fine.
(b)  Clearly not.  For this to work, we would need
$5n^2-3n+20<2n^2$ to hold as $n$ increases towards $\infty$.
But this would imply that $3n^2-3n+20<0$.
But when $n\geq 1$, $3n^2$ is positive and larger than $3n$, so
$3n^2-3n+20>0$.
(c) Sure.  The proof used the fact that the inequality is true when $n\geq 1$,
so it is clearly also true if $n\geq 100$.  And the definition of Big-O does
not require that we use the smallest possible value for $n_0$.
(d) No.  We would need a constant $c$ such that $5\cdot 0^2-3\cdot 0 + 20 = 20\leq 0= c\cdot 0^2$,
which is clearly impossible.
\end{Answer}
\begin{Answer}{6.9}
If $n\geq 1$,
\begin{eqnarray*}
5n^{5} -4n^4 + 3n^3 - 2n^2 +n
&\leq&5n^{5} + 3n^3 +n\\
&\leq&5n^{5} + 3n^5 +n^5\\
&=&9n^5.
\end{eqnarray*}
Therefore, $5n^{5} -4n^4 + 3n^3 - 2n^2 +n=O(n^5)$.
\end{Answer}
\begin{Answer}{6.10}
We used $n_0=1$ and $c=9$.
Your values for $n_0$ and $c$ may differ.
This is O.K. if you have the correct algebra to back it up.
\end{Answer}
\begin{Answer}{6.13}
Since $4n^2\leq 4n^2+n+1$ for $n\geq 0$, $4n^2=\Omega(n^2)$.
\end{Answer}
\begin{Answer}{6.14}
We used $c=4$ and $n_0=0$.  You might have used $n_0=1$ or some other positive value.
As long as you chose a positive value for $n_0$, it works just fine.  You could have also
used any value for $c$ larger than $0$ and at most $4$.
\end{Answer}
\begin{Answer}{6.17}
It is O.K.  Since the second inequality holds when $n\geq 0$, it also holds when $n\geq 1$.

In general, when you want to combine inequalities that contain two different assumptions,
you simply make the more restrictive assumption.  In this case, $n\geq 1$ is more restrictive
than $n\geq 0$.  In general, if you have assumptions $n\geq a$ and $n\geq b$, then to combine
the results with these assumptions, you assume $n\geq max(a,b)$.
\end{Answer}
\begin{Answer}{6.22}
$g(n)$ appears in the denominator of a fraction.  If at some point it does not become
(and remain) non-zero, the limit in the definition will be undefined.
If you never took a calculus course and are not that familiar with limits,
do not worry a whole lot about this subtle point.
\end{Answer}
\begin{Answer}{6.23}
$o$ is like $<$ and $\omega$ is like $>$.
\end{Answer}
\begin{Answer}{6.24}
(a) No.  If $f(n)=\Theta(g(n))$, $f$ and $g$ grow at the same rate.  But
$f(n)=o(g(n))$ expresses the idea that $f$ grows slower than $g$.
It is impossible for $f$ to grow at the same rate as $g$ and slower than $g$.
(b) Yes!  If $f$ grows no faster than $g$, then it is possible that it grows slower.
For instance, $n=O(n^2)$ and $n=o(n^2)$ are both true.
(c) No.  If $f$ and $g$ grow at the same rate, then $f(n)=O(g(n))$, but
$f(n)\neq o(g(n))$. For instance, $3n=O(n)$, but $3n\neq o(n)$.
(d) Yes.  In fact, it is guaranteed!  If $f$ grows slower than $g$, then
$f$ grows no faster than $g$.
%(e) No.  It is possible, but not certain.  For instance,
%$n^2=\Omega(n)$, but $n^2\neq\Theta(n)$.
\end{Answer}
\begin{Answer}{6.25}
\begin{solevalenum}
\item
Although this proof sounds somewhat reasonable, it is way too informal and convoluted.  Here are
some of the problems.
\begin{enumerate}
\item
This student misunderstands the concept behind `ignoring the constants.'
We can ignore the constants {\em after we know that $f(n)=O(g(n))$}.  We can't ignore them
in order to prove it.
\item
The phrase `become irrelevant' (used twice) is not precise.
We have developed mathematical notation for
a reason---it allows us to make statements like these precise.
It's kind of like saying that a certain car costs `a lot'.  What is `a lot'?
Although \$30,000 might be a lot for most of us, people with a lot more money
than I have might not think that \$500,000 is a lot.
\item
The phrase `This leaves us with $n^k+n^{k-1}+\cdots+n=O(n^k)$' is odd.  What precisely do they mean?
That this is true or that this is what we need to prove now?  In either case, it is incorrect.
Similarly for the second time they use the phrase `This leaves us with'.
\item
The second half of the proof is unnecessarily convoluted.  They essentially are claiming that their
proof has boiled down to showing that $n^k=O(n^k)$.
To prove this, they use an incredibly drawn out, yet vague, explanation that is
in a single unnecessarily long sentence.  Why are they even bringing $\Theta$ and $\Omega$ into this proof?
Why don't they just say something like `since $n^k\leq 1n^k$ for all $n\geq 1$, $n^k=O(n^k)$'?
I believe the answer is obvious:  they don't really understand what they are doing here.
They clearly have a vague understanding of the notation, but they don't understand the
formal definition.
\end{enumerate}
The bottom line is that this student understands that the statement they needed to prove is correct,
and they have a vague sense of {\em why} it is true, but they did not have a clear understanding
of how to use the definition of Big-O to prove it.
The most important thing to take away from this example is this:
{\em Be precise, use the notation and
definitions you have learned, and if your proofs look a lot different than those in the book, you
might question whether or not you are on the right track.}
\item
This proof is correct.
\end{solevalenum}
\end{Answer}
\begin{Answer}{6.26}
We cannot say anything about the relative growth rates of $f(n)$ and $g(n)$
because we are only given upper bounds for each.
It is possible that $f(n)=n^2$ and $g(n)=n$, so that
$f(n)$ grows faster, or vice-versa.  They could also both be $n$.
\end{Answer}
\begin{Answer}{6.27}
(a) F. This is saying that $f(n)$ grows {\em no} faster than $g(n)$.
(b) F. They grow at the same rate.
(c) F. $f(n)$ might grow slower than $g(n)$.  For instance, $f(n)=n$ and $g(n)=n^2$.
(d) F. They might grow at the same rate.  For instance, $f(n)=g(n)=n$.
(e) F. If $f(n)=n$ and $g(n)=n^2$, $f(n)=O(g(n))$, but $f(n)\neq \Omega(g(n))$.
(f) T. By Theorem~\ref{thm:theta}.
(g) F. If $f(n)=n$ and $g(n)=n^2$, $f(n)=O(g(n))$, but $f(n)\neq \Theta(g(n))$.
(h) F. If $f(n)=n$ and $g(n)=n^2$, $f(n)=O(g(n))$, but $g(n)\neq O(f(n))$.
\end{Answer}
\begin{Answer}{6.37}
$c_1 g(n)\leq f(n) \leq c_2 g(n) \mbox{ for all } n\geq n_0$;
$\frac{1}{ c_2} f(n)$;
$\frac{1 }{ c_1} f(n)$;
$c_3 h(n)\leq g(n) \leq c_4 h(n) \mbox{ for all } n\geq n_1.$;
$c_2$; $c_2 c_4$;
$\max\{n_0,n_1\}$;
$c_1 c_3 h(n)$;
$c_2 c_4 h(n)$;
$\Theta(h(n))$;
$\Theta$;
transitive;
\end{Answer}
\begin{Answer}{6.39}
(a) T. By Theorem~\ref{thm:asyprops4}.
(b) T. By Theorem~\ref{thm:theta}.
(c) T. By Theorem~\ref{thm:asyprops3}.
(d) F. The backwards implication is true, but the forward one is not.  For instance,
if $f(n)=n$ and $g(n)=n^2$, clearly $f(n)=O(g(n))$, but $f(n)\neq \Theta(g(n))$.
(e) F. Neither direction is true.  For instance,
if $f(n)=n$ and $g(n)=n^2$, $f(n)=O(g(n))$, but $g(n)\neq O(f(n))$.
(f) T. By Theorem~\ref{thm:asyprops4}.
(g) T. By Theorem~\ref{thm:theta}.
(h) T. By Theorem~\ref{thm:asyprops}.
\end{Answer}
\begin{Answer}{6.42}
$c_{1}n^2$;
$c_{2}n^2$;
$\frac{1}{2} - \frac{3}{n}$;
$\frac{10-6}{20}=\frac{1}{5}$;
$\frac{1}{5}n^2$;
$\frac{1}{2}n^2$;
 $10$.
\end{Answer}
\begin{Answer}{6.43}
There are a few ways to think about this.  First, the larger $n$ is, the smaller $\frac{3}{n}$ is,
so a smaller amount is being subtracted.  But that's perhaps too fuzzy.  Let's look at it this way:
$$n\geq 10\Rightarrow
 \frac{10}{3} \leq \frac{n}{3}
\Rightarrow \frac{3}{n}\leq \frac{3}{10}
\Rightarrow -\frac{3}{n}\geq -\frac{3}{10}
\Rightarrow \frac{1}{2}-\frac{3}{n}\geq \frac{1}{2}-\frac{3}{10}.
$$
\end{Answer}
\begin{Answer}{6.45}
(a) Theorem~\ref{thm:theta}.
(b) Absolutely not!  Theorem~\ref{thm:theta} requires that we also prove
$f(n)=\Omega(g(n))$.  Here is a counterexample: $n=O(n^2)$, but $n\neq\Theta(n^2)$.
So $f(n)=O(g(n))$ does not imply that $f(n)=\Theta(g(n))$.
\end{Answer}
\begin{Answer}{6.46}
Notice that when $n\geq 1$,
$n!=1\cdot 2\cdot 3\cdots n \leq n\cdot n\cdots n = n^n.$
Therefore $n!=O(n^n)$ (We used $n_0=1$, and $c=1$.)
\end{Answer}
\begin{Answer}{6.49}
If $f(x)=O(g(x))$,
then there are positive constants $c_1$ and $n_{0}^{\prime}$
such that
$$0\leq f(n)\leq c_1\,g(n)\mbox{ for all }n \ge n_{0}^{\prime},$$
and if $g(x)=O(h(x))$,
then there are positive constants $c_2$ and
$n_{0}^{\prime\prime}$ such that
$$0\leq g(n)\leq c_2\,h(n)\mbox{ for all }n \ge n_{0}^{\prime\prime}.$$
Set $n_0=\max(n_0^{\prime},n_0^{\prime\prime})$ and $c_3=c_1\,c_2$.
Then
$$0\leq f(n)\leq c_1\,g(n)\leq c_1\,c_2\,h(n)=c_3\,h(n)
\mbox{ for all }n \ge n_{0}.$$
Thus $f(x)=O(h(x))$.
\end{Answer}
\begin{Answer}{6.53}
(a) $\infty$
(b) $\infty$
(c) $\infty$
(d) $\infty$
(e) $0$
(f) $0$
(g) $8675309$
\end{Answer}
\begin{Answer}{6.57}
Theorem~\ref{limitProps1} part (b) implies that
$\displaystyle\lim_{n\rightarrow\infty} n^2=\infty$.
Since the limit being computed was actually
$\displaystyle \lim_{n\rightarrow\infty}\frac{1}{n^2}$,
Theorem~\ref{thm:oneoverlimit}
was used to obtain the final answer of $0$ for the limit.
\end{Answer}
\begin{Answer}{6.58}
Notice that $\lim_{x\rightarrow\infty}\frac{3x^3}{x^2}=
\lim_{x\rightarrow\infty}3x=\infty,$
so $3x^3=\omega(x^2)$ by the second case of the Theorem~\ref{thm:limits}, which also implies
that $3x^3=\Omega(x^2)$.
\end{Answer}
\begin{Answer}{6.63}
Notice that
${\displaystyle \lim_{n\rightarrow\infty}\frac{n(n+1)/2}{n^2} =
\lim_{n\rightarrow\infty}\frac{n^2+n}{2n^2} =
\lim_{n\rightarrow\infty}\frac{1}{2} +\frac{1}{2n}= \frac{1}{2}+0=\frac{1}{2}}$,
so $n(n+1)/2=\Theta(n^2)$.
\end{Answer}
\begin{Answer}{6.64}
(a)
Since
${\displaystyle \lim_{x\rightarrow\infty} \frac{2^x}{ 3^x}
= \lim_{x\rightarrow\infty} \left(\frac{2}{ 3}\right)^x
= 0}$, the result follows.

(b)
If $x\geq 1$, then clearly $(3/2)^x\geq 1$, so
$2^x\leq
2^x\left(\frac{3}{ 2}\right)^x=
\left(\frac{2\times 3}{ 2}\right)^x=3^x.$
Therefore, $2^x=O(3^x)$.
\end{Answer}
\begin{Answer}{6.70}
\begin{pfevalenum}
\item
$7^x$ grows faster than $5^x$ does not mean $7^x-5^x>0$ for all $x\neq 0$.
For one thing,
we are really only concerned about positive values of $x$.
Further, we are specifically concerned about very large values of $x$.
In other words, we want something to be true for all $x$ that are `large enough'.
Also, this statement does not take into account constant factors.
Similarly, a tight bound does {\em not} imply that $7^x-5^x=0$.
The bottom line:  This one is way off.  They are not conveying an understanding of
what `upper bound' really means, and they certainly haven't proven anything.
Frankly, I don't think they have a clue what they are trying to say in this proof.
\item
This one has several problems.  First, the application of l'Hopital's rule is incorrect.
The result should be $\displaystyle\lim_{x\rightarrow \infty}\frac{5^x\log 5}{7^x\log 7}$,
which should make it obvious that l'Hopital's rule doesn't actually help in this case.
(The key to this one is to do a little algebra.)
The next problem is the statement `but $x\log 7$ gets there faster'.
What exactly does that mean?
Asymptotically faster, or just faster?  If the former, it needs to be proven.  If the latter, that isn't enough to prove relative growth rates.  Finally, even if this showed that $5^x=O(7^x)$, that only shows
that $7^x$ is an upper bound on $5^x$.  It does not show that the bound is not tight.
The bottom line is that bad algebra combined with vague statements falls way short of
a correct proof.
\item
This proof is {\em very} close to being correct.  The main problem is that they only stated
that $5^x=O(7^x)$, but they also needed to show that $5^x\neq\Theta(7^x)$.  It turns out that the
theorem they mention also gives them that. So all they needed to add is `and $5^x\neq \Theta(7^x)$' at the end.
Technically, there is another problem---they should have taken the limit of $5^x/7^x$.  What they really showed
using the limit theorem is that $7^x=\omega(5^x)$, which is equivalent to $5^x=o(7^x)$.  It isn't a major problem,
but technically the limit theorem does not directly give them the result they say it does.
If you are trying to prove that $f(x)$ is bounded by $g(x)$, put $f(x)$ on the top and $g(x)$ on the bottom.
\end{pfevalenum}
\end{Answer}
\begin{Answer}{6.72}
You should have come up with $n^2\log n$ for the upper bound.  If you didn't,
now that you know the answer, go back and try to write the proofs before
reading them here.
(a)
If $n> 1$,
$$\ln (n^2+1) \leq \ln (n^2+n^2)= \ln (2 n^2) =
( \ln 2+ \ln n^2) \leq ( \ln n+ 2\ln n) =3\ln n$$
Thus when $n > 1$,
$$ n\ln (n^2+1) + n^2\ln n\leq n 3\ln n+n^2\ln n\leq
3 n^2 \ln n+n^2\ln n\leq  4 n^2 \ln n. $$
Thus, $n\ln (n^2+1) + n^2\ln n=O(n^2 \ln n).$
(You may have different algebra in your proof.  Just make certain
that however you did it that it is correct.)

(b)
$
\begin{array}[t]{llll}
\displaystyle\lim_{x\rightarrow\infty} \dfrac{{n\ln (n^2+1) + n^2\ln n}}{ n^2\ln n}
& =& \displaystyle\lim_{x\rightarrow\infty}
\dfrac{n\ln (n^2+1)}{ n^2\ln n}  + 1\\[.125in]
& =& \displaystyle1+ \lim_{x\rightarrow\infty}
\dfrac{\ln (n^2+1)}{ n\ln n} \\[.125in]
& =& \displaystyle1+ \lim_{x\rightarrow\infty}
\dfrac{\dfrac{2n}{n^2+1}}{1\cdot\ln n+n\cdot\frac{1}{n}} &\mbox{(l'Hopital)}\\[.125in]
& =& \displaystyle1+ \lim_{x\rightarrow\infty}
\dfrac{2n}{(n^2+1)(\ln n+1)} \\[.125in]
& =&\displaystyle 1+ \lim_{x\rightarrow\infty}
\dfrac{2}{2n(\ln n + 1) + (n^2+1)\cdot\frac{1}{n}} &\mbox{(l'Hopital)}\\[.125in]
& =&\displaystyle 1+ \lim_{x\rightarrow\infty}
\dfrac{2}{2n(\ln n + 1) + n+\frac{1}{n}}\\
&=&1+0=1.
\end{array}
$\\
Therefore,  $n\ln (n^2+1) + n^2\ln n=\Theta(n^2\log n)$.
\end{Answer}
\begin{Answer}{6.74}
We can see that $(n^2-1)^5=\Theta(n^{10})$ since
$$\lim_{n\rightarrow\infty}\frac{(n^2-1)^5}{n^{10}}
=\lim_{n\rightarrow\infty}\left(\frac{n^2-1}{n^2}\right)^5
=\lim_{n\rightarrow\infty}\left(1 - \frac{1}{n^2}\right)^5
=1
.$$
\end{Answer}
\begin{Answer}{6.75}
The following limit shows that $2^{n+1}+5^{n-1}=\Theta(5^n)$.
$$\lim_{n\rightarrow\infty}\frac{2^{n+1}+5^{n-1}}{5^n}
=\lim_{n\rightarrow\infty}\frac{2^{n+1}}{5^n}+\frac{5^{n-1}}{5^n}
=\lim_{n\rightarrow\infty}2\left(\frac{2}{5}\right)^n+\frac{1}{5}
=0+\frac{1}{5}.
$$
Note that we could also have shown that $2^{n+1}+5^{n-1}=\Theta(5^{n-1})$,
but that is not as simple of a function.
\end{Answer}
\begin{Answer}{6.78}
Since $a<b$, $b-a>0$.
Therefore, $\displaystyle \lim_{n\rightarrow\infty} \frac{n^a}{n^b}
=  \lim_{n\rightarrow\infty} n^{a-b} =  \lim_{n\rightarrow\infty} \frac{1}{n^{b-a}}=0$.
By Theorem~\ref{thm:limits}, $n^a=o(n^b)$.
\end{Answer}
\begin{Answer}{6.81}
Since $a<b$, $a/b<1$.
Therefore, $\displaystyle \lim_{n\rightarrow\infty} \frac{a^n}{b^n}
=  \lim_{n\rightarrow\infty}\left(\frac{a}{b}\right)^n = 0$.
By Theorem~\ref{thm:limits}, $a^n=o(b^n)$.
\end{Answer}
\begin{Answer}{6.86}
(a) False since $3^n$ grows faster than $2^n$.
(b) True since $2^n$ grows slower than $3^n$.
(c) False since $3^n$ grows faster than $2^n$, which means it does not grow
slower or at the same rate.
(d) True since they both have the same growth rate.  Remember, exponentials with
different bases have different growth rates, but logarithms with different bases
have the same growth rate.
(e) True since they have the same growth rate.  Remember that if $f(n)=\Theta(g(n))$,
then $f(n)=O(g(n))$ and $f(n)=\Omega(g(n))$.
(f) False since they have the same growth rate, so $\log_{10} n$ does not grow slower
than $\log_3 n$.
\end{Answer}
\begin{Answer}{6.89}
Using l'Hopital's rule, we have
$\displaystyle \lim_{n\rightarrow\infty} \frac{\log_c(n)}{n^b}
=  \lim_{n\rightarrow\infty}\frac{\frac{1}{n\ln(c)}}{b\, n^{b-1}}
=  \lim_{n\rightarrow\infty}\frac{1}{\ln(c)b\, n^b}
=0$ since $b>0$.  Thus, Theorem~\ref{thm:limits} tells us that $\log_c n=o(n^b)$.
\end{Answer}
\begin{Answer}{6.94}
(a) $\Theta$;
(b) $o$ ($O$ is correct, but not precise enough.);
(c) $\Theta$;
(d) $o$ ($O$ is correct, but not precise enough.);
(e) $\Theta$ since $2^n=2\,2^{n-1}$;
(f) $\Omega$ (Technically it is $\omega$, but I'll let it slide if you put $\Omega$ since
 we haven't used $\omega$ much.);
(g) through (j) are all $o$ ($O$ is correct, but not precise enough.)
\end{Answer}
\begin{Answer}{6.96}
If your answers do not all start with $\Theta$, go back and redo them before reading the answers.  Your answers should match the following {\em exactly}.
(a) $\Theta(n^7)$.
(b) $\Theta(n^8)$.
(c) $\Theta(n^2)$.
(d) $\Theta(3^n)$.
(e) $\Theta(2^n)$.
(f) $\Theta(n^2)$.
(g) $\Theta(n^{.000001})$.
(h) $\Theta(n^n)$.
\end{Answer}
\begin{Answer}{6.97}
Here is the correct ranking (where $\sim$ indicates two functions grow at the same rate):

\noindent
$10000$,\,
$\log x$ $\sim$ $\log (x^{300})$,\,
$\log^{300} x$,\,
$x^{.000001}$, \,
$x$ $\sim$ $\log(2^x)$,  \,
$x \log(x)$,  \,
$x^{log_2 3}$, \,
$x^2$, \,
$x^{5}$, \,
$2^x$,  \,
$3^x$.
\end{Answer}
\begin{Answer}{6.98}
Modern computers use multitasking to perform several tasks (seemingly) at the same time.
Therefore, if an algorithm takes 1 minute of real time (wall-clock time), it might be that
58 seconds of that time was spent running the algorithm, but it could also be the case that
only 30 seconds of that time were spent on that algorithm, and the other 30 seconds spent
on other processes.  In this case, the CPU time would be 30 seconds, but the wall-clock time
60 seconds.

Further complicating matters is increasing availability of machines
with multiple processors.  If an algorithm runs on 4 processors rather than one, it
might take 1/4th the time in terms of wall-clock time, but it will probably take the
same amount of CPU time (or close to it).
\end{Answer}
\begin{Answer}{6.99}
We cannot be certain whose algorithm is better with the given information.
Maybe Sue used a {\em TRS-80 Model IV} from the 1980s to run her program and Stu used
{\em Tianhe-2} (The fastest computer in the world from about 2013-2014).
In this case, it is possible that if Sue ran her program on {\em Tianhe-2}
it would have only taken 2 minutes, making her the real winner.
\end{Answer}
\begin{Answer}{6.100}
As has already been mentioned, other processes on the machine can have a significant influence
on the wall-clock time.  For instance, if I run two CPU-intensive programs at once, the wall-clock time of each might be about twice what it would be if I ran them one at a time.
If they are run on a machine with multiple cores
the wall-clock time might be closer to the CPU-time.  But other processes that are
running can still throw off the numbers.
\end{Answer}
\begin{Answer}{6.101}
For the most part, yes.  This is especially true if the running times of the algorithms
are not too close to each other (in other words, if one of the algorithms is significantly
faster than the other).
However, the number of other processes running on the machine
can have an influence on CPU-time.  For instance, if there are more processes running,
there are more context switches, and depending on how the CPU-time is counted, these context
switches can influence the runtime.  So although comparing the CPU-time of two algorithms
that are run on the same computer gives a pretty good indication of which is better, it
is still not perfect.
\end{Answer}
\begin{Answer}{6.103}
This one is a little more tricky.
The answer is $n\cdot m$ since this is how many entries are in the matrix.
Sometimes we need to use two numbers to specify the input size.
As suggested previously, we will ignore the size of the two other pieces of data.
\end{Answer}
\begin{Answer}{6.109}
We focus on the assignment (=) inside the loop and ignore the other instructions.
This should be fine since assignment occurs at least as often as any other instruction.
In addition, it is important to note that \verb+max+ takes constant time
(did you remember to explicitly say this?),
as do all of the other operations, so we aren't under-counting.
It isn't too difficult to see that the assignment will occur $n$ times for an array of size $n$
since the code goes through a loop with $i=0,\ldots, n-1$.
Thus, the complexity of \scode{maximum} is always $\Theta(n)$.  That is, $\Theta(n)$ is the best, average, and worst-case complexity of \scode{maximum}.
\end{Answer}
\begin{Answer}{6.112}
The line in the inner for loop takes constant time (let's call it $c$).
The inner loop executes $k=50$ times, each time doing $c$ operations.
Thus the inner loop does $50\cdot c$ operations, which is still just a constant.
The outer loop executes $n$ times, each time executing the inner loop, which takes $50\cdot c$ operations.
Thus, the whole algorithm takes $50\cdot c\cdot n=\Theta(n)$ time.
\end{Answer}
\begin{Answer}{6.113}
The line in the inner for loop takes constant time (let's call it $c$).
The inner loop executes $n^2$ times since $j$ is going from
$0$ to $n^2-1$, so each time the inner loop executes, it does $cn^2$ operations.
The outer loop executes $n$ times, each time executing the inner loop.
Thus, the total time is $n\times cn^2=\Theta(n^3)$.

This is an example of an algorithm with a double-nested loop
that is {\em worse than} $\Theta(n^2)$.  The point of this exercise is to make
it clear that you should never jump to conclusions too quickly when analyzing
algorithms.  Read the limits on loops very carefully!
\end{Answer}
\begin{Answer}{6.116}
(a) \verb+AreaTrapezoid+ is constant.
(b) \verb+factorial+is not constant.
It should be easy to see that it has a complexity of $\Theta(n)$.
(c)  \verb+absoluteValue+ is constant if we assume \verb+sqrt+ takes constant time.
\end{Answer}
\begin{Answer}{6.121}
Although it has a nested loop, the inside loop always executes $6$ times, which
is a constant.  So the algorithm takes about $6\cdot c\cdot n=\Theta(n)$ operations, not $\Theta(n^2)$.
\end{Answer}
\begin{Answer}{6.126}
(a) Since \verb+factorial+ has a complexity of $\Theta(n)$, it is not quadratic.
(b) Since there are $n^2$ entries it to consider,
the algorithm takes $\Theta(n^2)$ time, so it would be quadratic.\footnote{Technically this is linear with respect to the {\em size} of the input since the
size of the input is $n^2$.  But it is quadratic in $n$.  In either case, it is
$\Theta(n^2)$.}
\end{Answer}
\begin{Answer}{6.127}
Bubble sort, selection sort, and insertion sort are three of them that you may have seen before.
\end{Answer}
\begin{Answer}{6.133}
As we mentioned in our analysis, executing the conditional statement takes about
$3$ operations, and if it is true, about $3$ additional operations are performed.  So
the worst case is no more than about twice as many operations as the best case.  In other
words, we are comparing $c\cdot n^2$ to $2c\cdot n^2$, both of which are $\Theta(n^2)$.
\end{Answer}
\begin{Answer}{6.136}
Since both of these methods require accessing the $i$th element of the list
for some integer $i$, and since we must traverse the list from the head, clearly
the complexity of both methods is $\Theta(i)$.

We could be less specific and say that the complexity is $\Theta(n)$ since $0\leq i <n$.
However, when analyzing algorithms that make repeated calls to these methods,
using $\Theta(i)$ might give a more accurate answer overall.  It makes the analysis
more difficult, but sometimes it is worth it.

{\em Note: For doubly-linked lists,
some implementations traverse starting at the tail if the index is closer to the end
of the list.  However, that just means the complexity is no worse than $\Theta(n/2)=\Theta(n)$.
In other words, it only changes the complexity by a constant factor.}
\end{Answer}
\begin{Answer}{6.138}
All of them should be $\Theta(1)$, assuming we keep track of how many elements
are currently in the stack (which is a reasonable thing to do).
\end{Answer}
\begin{Answer}{6.139}
For an array, either enqueue or dequeue (but not both) will be $\Theta(n)$.
All of the others will be $\Theta(1)$, assuming we keep track of how many elements
are currently in the queue.  Note that the advantage of the circular array implementation
is that both enqueue and dequeue are $\Theta(1)$.
\end{Answer}
\begin{Answer}{6.140}
For the array implementation, addToFront, removeFirst and contains will all be $\Theta(n)$
and the rest will be $\Theta(1)$.
For the linked list implementation, contains will be $\Theta(n)$ and the rest will be
constant if we assume there is both a head and tail pointer.  If there is no tail pointer,
addToEnd will be $\Theta(n)$.
\end{Answer}
\begin{Answer}{6.141}
For unbalanced, all operations can be done in $\Theta(h)$ time, where $h$ is the height
of the tree.  This is the best answer you can give.  You could also say $O(n)$ time since
the height is no more then $n$, but this answer is not precise enough to be of much use.
You {\em cannot} say $\Theta(\log n)$ since
this is not necessarily true for an unbalanced tree.

For balanced (red-black, AVL, etc.), all operations can be implemented with complexity $\Theta(\log n)$.
\end{Answer}
\begin{Answer}{6.142}
The average-case complexity for all of these operations is $\Theta(1)$,
and the worst-case complexity is $\Theta(n)$.
\end{Answer}
\begin{Answer}{6.143}
\begin{solevalenum}
\item
I have no idea what logic they are trying to use here.  Sure, $a^n$ is an exponential
function, but what does that have to do with how long this algorithm takes?
This solution is way off.
\item
Having the $i$ in the answer is nonsense since it doesn't mean
anything in the context of a complexity---it is just a variable that happens to be
used to index a loop.  Further, the answer should be given using
$\Theta$-notation.  So this solution is just plain wrong.
Since having an $i$ in the complexity does not make any sense, this person either
has a fundamental misunderstanding of how to analyze algorithms or they didn't
think about their final answer.  Don't be like this person!
\item
This solution is O.K., but it has a slight problem.  Although the analysis given estimates
the worst-case behavior, it over-estimates it.  By replacing $i$ with $n-1$, they are
over-estimating how long the algorithm takes.
The call to \verb+pow+ only takes $n-1$ time once.
This solution can really only tell us
that the complexity is $O(n^2)$.  Is it possible the over-estimation of time resulted
in a bound that isn't tight?  Even if it turns out that $\Theta(n^2)$ is the correct bound,
this solution does not prove it.
Although they are on the right track,
this person needed to be a little more careful in their analysis.
\end{solevalenum}
\end{Answer}
\begin{Answer}{6.144}
Before you read too far:
if you did not use a summation in your solution, go back and try again!
This is very similar to the analysis of \verb+bubblesort+.  The for loop takes $i$ from
$0$ to $n-1$, and each time the code in the loop takes $i$ time (since that is how long
\verb+power(a,i)+ takes).  Thus, the complexity is
$$\sum_{i=0}^{n-1} i = \frac{(n-1)n}{2}=\Theta(n^2).$$
Notice that just because the answer is $\Theta(n^2)$, that does not mean that the
third solution to Evaluate~\ref{eval:addpow} was correct.  As we stated in the solution
to that problem, because they overestimated the number of operations,
they only proved that the algorithm has complexity $O(n^2)$.
\end{Answer}
\begin{Answer}{6.145}
Here is one solution.
\begin{code}
	  double addPowers(double a, int n) {
	      if(a==1) {
	          return n;
	      } else {
	          double sum = 1; // for the $a^0$ term.
	          double pow = 1;
	          for(int i=1;i<n;i++) {
	              pow = pow*a;
	              sum += pow;
	          }
	          return sum;
	      }
	  }
\end{code}
If $a=1$, the algorithm takes constant time.
Otherwise, it executes a constant number of operations and a single for loop $n$ times.
The code in the loop takes constant time.
Thus the algorithm takes $\Theta(n)$ time.
\end{Answer}
\begin{Answer}{6.146}
If you used recursion instead of a loop, cool idea.  However, go back and do it again.
There is an even {\em simpler} way to do it.  Need a hint?  Apply some of that
discrete mathematics material you have been learning!  When you have a solution that
does not use a loop or recursion (or you get stuck), keep reading.

The trick is to use the formula for a
geometric series (did you recognize that this is what \verb+addPowers+ is really computing?).
We need a special case for $a=1$ because the
formula requires that $a\neq 1$.

\begin{code}
	  double addPowers(double a, int n) {
	      if(a==1) {
	          return n;
	      } else {
	          return (1-power(a,n))/(1-a);
	      }
	  }
\end{code}

If $a=1$, the algorithm takes constant time.
Otherwise, it executes a constant number of operations and a single call
to \verb@power(a,n)@ which takes $n$ time.
Thus the algorithm takes $\Theta(n+1)=\Theta(n)$ time.

It is worth noting that $a=0$ is a tricky case.
\verb+addPowers+ can't really be computed for $a=0$ since $0^0$ is undefined.
It is for this reason that the first term of a geometric sequence is technically $1$, not $a^0$.
Since $a^0=1$ for all other values of $a$, the case of $a=0$ is usually glossed over.
If you don't understand what the fuss is about, don't worry too much about it.
\end{Answer}
\begin{Answer}{6.148}
\begin{solevalenum}
\item
This takes about $2+4n+4m$ operations, which is essentially the
same as the `C' version.  Unfortunately, it is slightly worse than the original
solution since it is now incorrect.  All they did is omit adding the final $n$ in the first
sum.  This went from a `C' to a `D' (at best).
\item
This one takes about $2+4(n-1-m+1)=2+4(n-m)$ operations.  This is a lot better than
the previous solutions.  Unfortunately, it misses adding the final $n$, so it is incorrect.
It also is not as efficient as possible.  I'd say this is still a `C'.
\item
This student figured out the trick---they know a formula to compute the sum,
so they tried to use it.
Unfortunately, they used the formula incorrectly and/or they made a mistake when
manipulating the sum (it is impossible to tell exactly what they did wrong---they
made either one or two errors), so the algorithm is not correct.
In terms of efficiency, their solution is great because it takes a constant
number of operations no matter what $n$ and $m$ are.
Because their answer is efficient and {\em very close} to being correct,
I'd probably give them a `B'.
\end{solevalenum}
\end{Answer}
\begin{Answer}{6.149}
We use the fact that
${\displaystyle \sum_{k=m}^n k = \sum_{k=1}^n - \sum_{k=1}^{m-1}
=\frac{n(n+1)}{2} - \frac{(m-1)m}{2} }$ to give us the following solution:
\begin{code}
    int sumFromMToN(int m,int n) {
	    return n*(n+1)/2 - (m-1)*m/2;
	}
\end{code}
Since this is just doing a fixed number of arithmetic operations no matter
what the values of $m$ and $n$ are, it takes constant time.
\end{Answer}
\begin{Answer}{6.152}
The analysis of these is very similar to the analysis of Examples~\ref{exa:retainAllAL}, so the details are omitted.
(a) For a LinkedList the \verb+contains+ method takes $\Theta(m)$ time, so the overall
complexity is $\Theta(nm)$.
(b)For a HashSet it takes $\Theta(1)$ time to call \verb+contains+ (on average), so
the overall complexity is $\Theta(n+n)=\Theta(n)$.
\end{Answer}
\begin{Answer}{6.154}
(a) contains takes $\Theta(m)$ time so the complexity is
$\Theta(n(\log n+m))$.

(b) Here the contains method takes $\Theta(\log m)$ time, so the overall
complexity is $\Theta(n(\log n+\log m))$
(or $\Theta(n\log n + n\log m)$
if you prefer to write it that way).

Note that we don't know the relationship between
$n$ and $m$ so we can't simplify either answer.
\end{Answer}
\begin{Answer}{6.157}
\begin{tabular}{|l|l|l|l|l|}
\hline
\multicolumn{2}{|c|}{$n$}&
\multicolumn{2}{c|}{$\lfloor n/2\rfloor$}\\
decimal&binary\, \, \, \, \, &decimal&binary\, \, \, \, \,\\
\hline
\hline
12&1100&6&110\\
\hline
13&1101&6&110\\
\hline
32&100000&16&10000\\
\hline
33&100001&16&10000\\
\hline
118&1110110&59&111011\\
\hline
119&1110111&59&111011\\
\hline
\end{tabular}
\end{Answer}
\begin{Answer}{6.158}
The next theorem answers the question about the pattern.
\end{Answer}
